\documentclass{situationist-preprint}

\title{Categorical Desires and Endless Life}
\msid{EXI-03-b}
\msfamily{EXI-03}
\msdiscipline{Meaning, Death and the Absurd}
\mstype{Consequence}
\mskeywords{immortality; categorical desires; Williams; boredom; death}

\begin{document}
\maketitle

\begin{abstract}
EXI-03-a argues that an endless life would eventually exhaust the categorical desires that give life its point, or cease to be a life that matters to the person now. We derive three consequences about categorical desires themselves. First, the desire to live forever is not a categorical desire that gives life its point, since its object is not a project. Second, the undesirability of endless life is compatible with the badness of death at every time at which actual people die, because exhaustion lies far beyond any actual life. Third, the result places no upper bound on the desirable length of a finite life: it shows only that some finite length would be enough. The argument thus constrains the wish for immortality without recommending an early death.
\end{abstract}

\section{The result and its terms}

A categorical desire is one not conditional on the person's being alive, whose satisfaction requires her continued life; it gives life its point if its object is a project she takes to be worthwhile for reasons independent of her desiring it (Definition~2.1 of EXI-03-a). Given a bounded repertoire of project types for any character (Assumption~2.3), the exhaustion of each project type pursued long enough (Assumption~2.4), and the dependence of what matters on connectedness (Assumption~2.5), Theorem~3.3 of EXI-03-a concludes that an endless life is not desirable for the one who would live it.

\section{Three consequences}

\begin{proposition}\label{p:wish}
The desire to live forever is not a categorical desire that gives life its point.
\end{proposition}

\begin{proof}[Argument]
A categorical desire gives life its point only if its object is a project the person takes to be worthwhile. The object of the desire to live forever is the continuation of life itself, without specification of what is to be done in it. Continuation is not a project: it is the condition of all projects, and its worth for the person is derivative from the worth of what would fill it. By Theorem~3.3 of EXI-03-a, an endless continuation would eventually contain nothing that gives life its point. The desire is thus parasitic on categorical desires that it would, if satisfied, outlast.
\end{proof}

This explains a familiar feature of the wish for immortality: it is vivid in prospect and empty in detail. Those who wish to live forever usually have in mind the continuation of their present projects, which are finite in number and each, by Assumption~2.4, exhaustible \citep{Williams1973}.

\begin{proposition}\label{p:death}
That an endless life is not desirable does not imply that death is good, or not bad, at any time at which actual people die.
\end{proposition}

\begin{proof}[Argument]
By Lemma~3.1 of EXI-03-a, exhaustion occurs only after some time in an endless life; the theorem does not say when. Every actual death occurs within a span in which people still have unexhausted categorical desires: projects unfinished and project types not yet taken up. At such a time the life the person would have had, had she not died, contains goods her actual life lacks, and by the deprivation account of EXI-02 her death is bad for her. The theorem concerns the limit of a life; the badness of death concerns its actual ends.
\end{proof}

This is the content of the remark in EXI-03-a that neither horn makes death good. The structure is the one \citet{Williams1973} himself adopted: it is a good thing that we are not immortal, and it can still be a misfortune to die when one does.

\begin{proposition}\label{p:length}
The result establishes that some finite length of life would suffice to exhaust the categorical desires a connected character can support. It does not establish what that length is, nor that it is near the present human span.
\end{proposition}

The repertoire of a character is finite, but it may be large, and a person's character develops within the bounds of connectedness and so opens new project types as others close. A life many times longer than ours might be desirable throughout. The argument therefore gives no support to the claim that a human lifespan is already long enough, and it is consistent with regret that it is not longer \citep{Fischer1994}.

\section{An objection}

\begin{objection}
If the desire for endless life is not categorical, then the fear of losing it should also be idle. Yet people who are told that they will not live forever are not relieved.
\end{objection}

\begin{reply}
Their concern is not for endlessness but for the categorical desires they have now, which death will frustrate whenever it comes before they are exhausted. Proposition~\ref{p:death} explains this concern, and Proposition~\ref{p:wish} explains why it attaches to particular projects and not to the bare prospect of continuing.
\end{reply}

\section{Conclusion}

Categorical desires give life its point and, being finite in an endless life, set a limit on the value of continuation. The wish to live forever outruns them; the wish to live longer than one will may not.

\begin{thebibliography}{9}\small
\bibitem[Fischer(1994)]{Fischer1994} Fischer, J.~M. (1994). Why immortality is not so bad. \emph{International Journal of Philosophical Studies}, 2(2), 257--270.
\bibitem[Williams(1973)]{Williams1973} Williams, B. (1973). The Makropulos case: Reflections on the tedium of immortality. In \emph{Problems of the Self}, 82--100. Cambridge University Press.
\end{thebibliography}

\end{document}
